% ===== 原文 pp.3–5 =====

% [原文 p.3]
$K$ 理论对 $X$ 的函子性如下。给定 $X\xrightarrow{f}Y$ 与向量丛 $E\xrightarrow{p}Y$，我们有拉回 $f^*E\xrightarrow{q}X$，
\[
f^*E=X\times_Y E=\{(x,e)\in X\times E: p(e)=f(x)\},\qquad q(x,e)=x,
\]
它是 $X$ 上的向量丛，纤维为 $q^{-1}(x)=p^{-1}(f(x))$。至于向量丛的各条条件是否满足，只需拉回 $E\xrightarrow{p}Y$ 的一个局部平凡化即可检验。容易验证，由此得到一个环同态 $K(Y)\xrightarrow{f^*}K(X)$。

事实上，$K$ 经由同伦范畴分解。

% [Tutor Note: p.3 「$K$ 经由同伦范畴分解」正是命题 2 的内容：同伦的 $f_0\simeq f_1$ 诱导同一拉回 $f_0^*=f_1^*$，故 $X\mapsto K(X)$ 可视为从同伦范畴出发的反变函子。这也说明命题 2 不是技术性附带结论：同伦不变性正是 $K$ 能成为广义上同调理论的前提。]

\begin{prop}{2}（[VBKT] 1.6，[AGP] 8.4.4.）若 $f_0,f_1:A\to B$ 经由 $H:A\times I\to B$ 同伦，且 $p:E\to B$ 是向量丛，则 $f_0^*(E)\cong f_1^*(E)$（假定 $A$ 紧 Hausdorff）。
\end{prop}

\begin{proof}
考虑拉回 $q:H^*(E)\to A\times I$，我们要证明
\[
f_0^*(E)=q^{-1}(A\times\{0\})\cong q^{-1}(A\times\{1\})=f_1^*(E).
\]
\begin{itemize}[leftmargin=1.6em,itemsep=2pt]
\item 首先我们注意：向量丛 $E\xrightarrow{p}X\times I$ 只要其在 $X\times[0,1/2]$、$X\times[1/2,1]$ 上的限制都平凡，它本身即平凡。若这两个限制为 $E_0,E_1$，并带有局部平凡化
\[
E_0\xrightarrow{h_0}X\times[0,1/2]\times\C^n,\qquad E_1\xrightarrow{h_1}X\times[1/2,1]\times\C^n,
\]
则我们有同构 $h_1^{-1}\circ h_0:X\times\{1/2\}\times\C^n\to X\times\{1/2\}\times\C^n$，它由 $(x,1/2,v)\mapsto(x,1/2,g(v))$ 给出，其中 $g$ 是某个线性映射；而映射
\[
X\times I\times\C^n\to B,\qquad (x,t,v)\mapsto
\begin{cases}
h_0(x,t,v) & 0\le t\le 1/2\\
h_1(x,t,g(v)) & 1/2\le t\le 1
\end{cases}
\]
是 $E$ 的一个平凡化。
% [Tutor Note: p.3 第一个 bullet 的展示式原文印作 $X\times I\times\C^n\to B$，但紧接着说它是「$E$ 的一个平凡化」，落点应为全空间 $E$（已用 700 dpi 裁图复核，原文确实印成 $B$，非排版误读）。另注意 $h_0,h_1$ 本是定义在 $E_0,E_1$ 上的局部平凡化，而分段式里 $h_0(x,t,v)$ 直接作用在平凡化坐标 $(x,t,v)$ 上，应理解为 $h_0^{-1}$；原文记法不严，正文照译。]
\item 其次我们注意：任一向量丛 $E\xrightarrow{p}B\times I$ 都有形如 $U_i\times I$ 的局部平凡化，其中 $U_i$ 是 $B$ 的开覆盖。给定 $x\in X$，存在 $x$ 的一族邻域 $\{U_{x,i}\}_{i=1}^n$ 以及一个分划 $0=t_0<\cdots<t_n=1$，使得 $E$ 在每个 $U_{x,i}\times[t_{i-1},t_i]$ 上的限制都平凡（这里用到 $I$ 的紧性）。由前一点，$E$ 在 $U_x\times I$ 上的限制平凡，其中 $U_x=\bigcap_{i=1}^n U_{x,i}$。
% [Tutor Note: p.3 $U_x=\bigcap_{i=1}^nU_{x,i}$ 是有限多个开集的交，故仍为开邻域；此处必须用到 $n$ 有限（由 $I$ 的紧性得到有限分割），无限交不保证开。核对方法：把 $I$ 换成非紧的 $(0,1)$ 或只取可数分割，上述论证即失效。]
\end{itemize}

现在取任一向量丛 $E\xrightarrow{p}X\times I$，并取 $X=\bigcup_{i=1}^n U_i$，使得 $E$ 在 $U_i$ 上平凡，且带有一个从属于诸 $U_i$ 的单位分解 $\varphi_i$。令 $\psi_i=\sum_{j=1}^i\varphi_j$，$X_i\subseteq X\times I$ 为 $\psi_i$ 的图像（同胚于 $X$），而 $E_i\xrightarrow{p_i}X_i$ 为 $E$ 在 $X_i$ 上的限制。由于 $E$ 在 $U_i\times I$ 上平凡，同胚 $h_i:X_i\to X_{i-1}$ 可提升为同胚 $\tilde h_i:E_i\to E_{i-1}$，且它们在 $p^{-1}(U_i\times I)$ 之外是恒等映射：
\[
\tilde h_i(x,\psi_i(x),v)=(x,\psi_{i-1}(x),v),\qquad (x,\psi_i(x),v)\in U_i\times I\times\C^n\cong p^{-1}(U_i\times I).
\]
复合 $\tilde h_1\circ\cdots\circ\tilde h_n$ 是同胚 $E|_{X\times\{0\}}\cong E|_{X\times\{1\}}$。
% [Tutor Note: p.3 注意符号复用：第一个 bullet 中的 $h_0,h_1$ 是局部平凡化，而本段的 $h_i:X_i\to X_{i-1}$ 是 $\psi_i$ 图像之间的同胚（把曲线 $X_i$ 沿 $t$ 方向「拉直」到 $X_{i-1}$），两者不是一回事。$\tilde h_i$ 是把 $h_i$ 逐纤维提升所得，故在 $p^{-1}(U_i\times I)$ 之外只能取恒等，这一点是核对「$\tilde h_i$ 是整体同胚」的关键。]
\end{proof}

一些例子：

\begin{itemize}[leftmargin=1.6em,itemsep=2pt]
\item $K(pt.)=\Z$（点上的向量丛就是一个向量空间，在同构意义下由它的维数完全决定）。

% [原文 p.4]
\item 我们将看到 $K(S^2)=\Z[H]/(H-1)^2$，它由典范丛 $H=\{(L,x):L\in P^1(\C),\,x\in L\}$ 生成，这里把 $S^2$ 看作 $\C^2$ 中过原点的全体复直线 $L$ 的集合。至少可以看出关系 $H\otimes H\oplus 1=H\oplus H$ 成立，所用的正是仿射开覆盖之交上的粘合函数
\[
f(z)=\begin{pmatrix} z^2 & 0 \\ 0 & 1\end{pmatrix},\qquad g(z)=\begin{pmatrix} z & 0 \\ 0 & z\end{pmatrix}.
\]
典范丛 $H$ 在两块平面 $[1:w]$、$[z:1]$ 上都平凡（$w=1/z$）。在它们的交集上，点 $(\lambda,\lambda w)\in[1:w]$ 经乘以 $z$ 映到点 $(\lambda z,\lambda)\in[z:1]$。因此 $H\otimes H\oplus 1$ 与 $H\oplus H$ 的粘合函数如上所述。映射 $\alpha_t:[0,1]\to \GL_2(\C)$
\[
\alpha_t=\begin{pmatrix}\cos(\pi t/2) & -\sin(\pi t/2)\\ \sin(\pi t/2) & \cos(\pi t/2)\end{pmatrix}
\]
是 $\GL_2(\C)$ 中连接
\[
\begin{pmatrix} 1 & 0 \\ 0 & 1\end{pmatrix}\qquad\text{与}\qquad\begin{pmatrix} 0 & 1 \\ 1 & 0\end{pmatrix}
\]
的一条道路，于是我们得到两个丛 $H\otimes H\oplus 1$、$H\oplus H$ 的粘合函数之间的同伦 $f\simeq g$，它由下式定义：
\[
\begin{pmatrix} z & 0 \\ 0 & 1\end{pmatrix}\alpha_t\begin{pmatrix} 1 & 0 \\ 0 & z\end{pmatrix}\alpha_t
\]
因此这两个丛同构，从而得到一个同态 $\Z[H]/(H-1)^2\to K(S^2)$。
% [Tutor Note: p.4 先核对 $f,g$ 的来历：$1$ 指一维平凡丛（转移函数恒为常数 $1$），$H$ 的转移函数是 $z$，故 $H\otimes H$ 的转移函数是 $z^2$、$H\otimes H\oplus 1$ 的是 $\mathrm{diag}(z^2,1)$，即 $f(z)$；而 $H\oplus H$ 的是 $\mathrm{diag}(z,z)$，即 $g(z)$。]
% [Tutor Note: p.4 末一个展示式需补算端点：$\alpha_0=I$，$\alpha_1=\begin{pmatrix}0&-1\\1&0\end{pmatrix}$，于是 $\begin{pmatrix}z&0\\0&1\end{pmatrix}\alpha_1\begin{pmatrix}1&0\\0&z\end{pmatrix}\alpha_1=-\begin{pmatrix}z^2&0\\0&1\end{pmatrix}=-f(z)$，并非恰好等于 $f(z)$。它与 $f(z)$ 只差左乘 $-I$，而 $-I$ 在 $GL_2(\C)$ 中经 $t\mapsto e^{i\pi t}I$ 与 $I$ 同伦，故同伦 $f\simeq g$ 的结论仍然成立。]
% [Tutor Note: p.4 本页有三处英文拼写笔误（正文照译，此处给出正确拼写，均已裁图复核）：①「via mulitiplication by $z$」中的 "mulitiplication"（1400 dpi）应为 "multiplication"；②「the cluthing for $H\otimes H\oplus 1$」（900 dpi）应为 "clutching"，即粘合函数；③定义 3 中「orthogoal complement」（1100 dpi）应为 "orthogonal complement"，即正交补。]
% [Tutor Note: p.4 定义 3 的两个展示式原文均印作 $v\in p$（小写 $p$，800 dpi 复核）；按上下文 $P$ 是 $n$-平面、$v$ 应是 $P$ 中的向量，正确形式为 $v\in P$。这是与上一条并列的原文笔误，正文照译以免「顺手改原文」。]
\end{itemize}

我们还有\textbf{约化 $K$ 理论}。在带点空间上，包含映射 $\{pt.\}\to X$ 诱导出映射 $K(X)\to\Z$，它把丛送到该丛在基点所在连通分支上的纤维的维数。$X$ 的约化 $K$ 理论 $\widetilde{K}(X)$ 就是这个环同态的核。一个基点的包含映射与到该基点的常值映射诱导出分裂 $K(X)\cong\widetilde{K}(X)\oplus\Z$。

\section{向量丛的分类}

\begin{defn}{3} \textbf{Grassmann 流形} $G_n(\C^N)$（$n\le N$）是 $\C^N$ 的全体 $n$ 维复向量子空间组成的集合，其拓扑取为 $U_N/(U_n\times U_{N-n})$（选定 $\C^N$ 的一组标准正交基之后，酉群在 $n$-平面上可迁地作用，稳定子群是那些固定该 $n$-平面及其正交补的矩阵）。在包含映射 $G_n(\C^N)\to G_n(\C^{N+1})$ 下取有向极限，我们得到 $G_n(\C^\infty)$。这些 Grassmann 流形带有\textbf{典范丛}
\[
E_n(\C^N)\to G_n(\C^N),\qquad E_n(\C^N)=\{(P,v):P\in G_n(\C^N),\,v\in p\},
\]
\[
E_n(\C^\infty)\to G_n(\C^\infty),\qquad E_n(\C^\infty)=\{(P,v):P\in G_n(\C^\infty),\,v\in p\}.
\]
\end{defn}

这些 $G_n(\C^N)$ 都是复维数为 $n(N-n)$ 的闭复流形。一个有用的表示是：把它看成 $n\times(N-n)$ 矩阵（其行张成该 $n$-平面）模去（左）$\GL_n$ 作用。于是每个 $n$-平面都有唯一的行简化阶梯形代表；按主元位置对这些代表分组，就给出 $G_n(\C^N)$ 上的一个胞腔结构。[关于特征映射的讨论，见 [VBKT] 命题 1.17。] 通过添加一列
% [原文 p.5]
零，我们得到子复形 $G_n(\C^N)\subseteq G_n(\C^{N+1})$。例如，$G_3(\C^5)$ 中的一个 4-胞腔嵌入 $G_3(\C^6)$ 之后可以呈这样的形状：
\[
\begin{pmatrix} 1 & * & 0 & 0 & * \\ 0 & 0 & 1 & 0 & * \\ 0 & 0 & 0 & 1 & * \end{pmatrix}
\mapsto
\begin{pmatrix} 0 & 1 & * & 0 & 0 & * \\ 0 & 0 & 0 & 1 & 0 & * \\ 0 & 0 & 0 & 0 & 1 & * \end{pmatrix}.
\]
$G_n(\C^\infty)$ 的元素可以表示为只有有限多个非零元的 $n\times\infty$ 行简化阶梯形矩阵。通过添加一整行零以及单独一个 $1$，我们得到子复形 $G_n(\C^\infty)\subseteq G_{n+1}(\C^\infty)$。例如，上面的 4-胞腔（先映入 $G_3(\C^\infty)$，再映入 $G_4(\C^\infty)$）成为
\[
\begin{pmatrix}
\dots & 0 & 1 & * & 0 & 0 & * & 0\\
\dots & 0 & 0 & 0 & 1 & 0 & * & 0\\
\dots & 0 & 0 & 0 & 0 & 1 & * & 0\\
\dots & 0 & 0 & 0 & 0 & 0 & 0 & 1
\end{pmatrix}.
\]
% [Tutor Note: p.5 例子里的矩阵是 $n\times N$（含主元列）的完整行简化阶梯形：$G_3(\C^5)$ 本该给 $3\times(5-3)=3\times2$ 矩阵，原图画的是 $3\times5$；右边 $G_3(\C^6)$ 的是 $3\times6$。核对方式：数一下左边矩阵 3 行 5 列、右边 3 行 6 列，恰是 $n\times N$ 而非 $n\times(N-n)$；正文「$n\times(N-n)$ 矩阵」应理解为「删去主元列后剩下 $n\times(N-n)$ 个自由元」。]
% [Tutor Note: p.5 大矩阵的读法（可逐格核对）：$\dots$ 左侧省略的全是零列，其后 7 列中主元依次在第 2、4、5、7 列；第一行第 7 列为 $0$ 而非 $*$，正因为新添的那个 $1$ 所在的列是新的主元列，而行简化阶梯形要求主元以上全为零。对照上一个小矩阵：$G_3(\C^6)$ 的矩阵（第一行 $[0,1,*,0,0,*]$）右端再添一列、并补上第 4 行 $(0,0,0,0,0,0,1)$，恰好得到所示矩阵。]

\begin{defn}{4} 有限与无限酉群的\textbf{分类空间}分别是
\[
BU_n=\cup_N G_n(\C^N)=G_n(\C^\infty),\qquad BU=\cup_n BU_n
\]
其中的映射来自上一段。
\end{defn}

我们借助下面这条命题，开始把 $K$ 理论与同伦、广义上同调联系起来。

\begin{prop}{3}（[VBKT] 1.16，[AGP] 8.5.13）从 $X$ 到 $\C^\infty$ 中 $n$-平面 Grassmann 流形的映射的同伦类，与 $X$ 上秩 $n$ 向量丛的同构类之间存在双射（假定 $X$ 紧 Hausdorff）：
\[
[X,BU_n]\to\Vect_\C^n(X),\qquad [f]\mapsto f^*(E_n(\C^\infty)).
\]
\end{prop}
% [Tutor Note: p.5 命题 3 陈述第 2 行的「ismophism」（900 dpi 裁图复核）应为 "isomorphism"，即同构；与 p.4 的三处拼写错误同属一批笔误，正文照译。]

\begin{proof}
设 $E\xrightarrow{p}X$ 是 $n$ 维向量丛。首先我们指出，同构 $E\cong f^*(E_n(\C^\infty))$ 等价于一个在每个纤维上都是线性单射的映射 $g:E\to\C^\infty$（[AGP] 把这样的 $g$ 称为 Gauss 映射）。若已有这样的同构，考虑下面的图
\begin{tikzcd}
E \arrow[r,"\cong"] \arrow[dr,"p"'] & f^*E_n(\C^\infty) \arrow[d] \arrow[r,"\tilde f"] & E_n(\C^\infty) \arrow[d] \arrow[r,"\pi"] & \C^\infty \\
& X \arrow[r,"f"] & BU_n &
\end{tikzcd}
（其中 $\pi(P,v)=v$）。复合 $g=\tilde f\circ\pi$ 在每个纤维上都是线性单射，因为 $\tilde f$ 与 $\pi$ 都是。反过来，给定这样的 $g$，令 $f(x)=g(p^{-1}(x))\in BU_n$，便得到同构 $E\cong f^*E_n(\C^\infty)$：
\[
v\in p^{-1}(x)\mapsto((f(x),g(v)),x)\in f^*E_n(\C^\infty).
\]
% [Tutor Note: p.5 命题 3 证明中「$g=\tilde f\circ\pi$」的复合顺序与图示相反：从 $E$ 出发沿上行依次是 $E\xrightarrow{\cong}f^*E_n(\C^\infty)\xrightarrow{\tilde f}E_n(\C^\infty)\xrightarrow{\pi}\C^\infty$，故应为 $g=\pi\circ\tilde f$（再与左端同构复合）；$\pi$ 定义在 $E_n(\C^\infty)$ 上，写成 $\tilde f\circ\pi$ 无法按字面复合。正文照译。]
\begin{itemize}[leftmargin=1.6em,itemsep=2pt]
\item {[满射性.]} 设 $E\xrightarrow{p}X$ 是 $n$ 维向量丛，$X=\bigcup_{i=1}^m U_i$ 是使 $E$ 在其上平凡的有限开覆盖，$\{\varphi_i\}_i$ 是从属于该覆盖的单位分解。对 $(x,v)\in p^{-1}(U_i)$，映射 $(x,v)\mapsto\varphi_i(x)v\in\C^n$ 可延拓为映射 $g_i:E\to\C^n$，它在 $p^{-1}(U_i)$ 之外为零。把它们拼在一起，便得到映射 $g:E\to(\C^n)^m\subseteq\C^\infty$，它在每个纤维上都是线性单射。
% [Tutor Note: p.5 本 bullet 的标签原文为 [Surjectivity.]（满射性），对应的另一方向标签是 p.6 开头的 [Injectivity.]（单射性）——两词是本节证明的两半，勿译成「单射」而混淆。另注意此处只用到有限覆盖 $m<\infty$，$\C^\infty$ 的「无限」体现在坐标的维数而非覆盖的个数。]
% [跨块环境：本文件止于 p.5 末句，itemize 与 proof 均有意不闭合；p.6 的 [Injectivity.]（单射性）bullet 与证明末的方框由 tr-02a 续写并收尾（见 BRIEF §4、Lead 裁决 2026-11-04）。请勿在本文件末尾补上 itemize/proof 的收尾命令。]
